% =============================================================================
% Journal manuscript template -- ASME Journal of Mechanical Design (JMD)
%
% Target venue : ASME Journal of Mechanical Design
% Class file   : asmejour.cls (CTAN: https://ctan.org/pkg/asmejour)
%                Auto-installed on demand by MiKTeX; provided by TeX Live.
% Documentation: https://mirrors.ctan.org/macros/latex/contrib/asmejour/asmejour-template.pdf
%
% This is a lorem-ipsum scaffold. Technical content has intentionally been
% omitted (see issue: "Manuscript venues and template"). See README.md for
% a summary of JMD author guidelines and an aside on Smart Materials and
% Structures (SMS) as a backup venue.
% =============================================================================

% PDF/A metadata (added in 2022; required by current asmejour for accessibility)
\DocumentMetadata{%
  lang        = en-US,
  pdfstandard = A-3u,
  pdfversion  = 1.7,
}

% Class options used here (see asmejour-template.pdf for the full list):
%   * lineno       -- numbered lines for review markup (run pdflatex twice)
%   * singlecolumn -- one column draft layout (drop for ASME two-column final)
%   * nocopyright  -- suppress ASME copyright footer until acceptance
%   * upint, varvw -- typographic preferences carried over from the template
%   * hyphenate    -- allow hyphenation in typewriter font
\documentclass[lineno,singlecolumn,nocopyright,upint,varvw,hyphenate]{asmejour}

\allowdisplaybreaks % from amsmath; multiline equations may break across pages

% --- Additional packages ---
% asmejour already loads: graphicx, hyperref, natbib, amsmath, amssymb,
% xcolor, booktabs, subcaption, caption, lineno, etc. Avoid re-loading those.
\usepackage{lipsum}      % lorem ipsum filler -- REMOVE before submission
\usepackage{siunitx}

% --- Draft helpers (remove for final submission) ---
\newcommand{\todo}[1]{\textcolor{red}{\textbf{[TODO: #1]}}}
\newcommand{\note}[1]{\textcolor{blue}{\textbf{[NOTE: #1]}}}

% --- PDF metadata ---
\hypersetup{%
  pdfauthor   = {Jeffrey R. Hill and Sterling G. Baird},
  pdftitle    = {Bayesian Optimization of Multi-Material 3D-Printed Tensegrity
                 Structures for Energy Absorption},
  pdfkeywords = {tensegrity, multi-material 3D printing, Bayesian optimization,
                 multifidelity, energy absorption},
  pdfsubject  = {Manuscript draft for ASME Journal of Mechanical Design},
}

% --- Journal name (printed in the running head; omit "Journal of") ---
\JourName{Mechanical Design}

\begin{document}

% =============================================================================
% Author block(s) -- one \SetAuthorBlock per author (or per shared affiliation)
% Mark corresponding author(s) with \CorrespondingAuthor (no space before).
% =============================================================================
\SetAuthorBlock{Jeffrey R. Hill\CorrespondingAuthor}{%
  Department of Mechanical Engineering,\\
  Brigham Young University,\\
  Provo, UT 84602, USA\\
  email: jeff.hill@byu.edu
}

\SetAuthorBlock{Sterling G. Baird}{%
  Department of Mechanical Engineering,\\
  Brigham Young University,\\
  Provo, UT 84602, USA\\
  email: sterling.baird@byu.edu
}

\title{Bayesian Optimization of Multi-Material 3D-Printed Tensegrity
       Structures for Energy Absorption}

% Keywords are printed at the end of the abstract; must precede \end{abstract}.
\keywords{tensegrity, multi-material 3D printing, Bayesian optimization,
          multifidelity, energy absorption}

% -----------------------------------------------------------------------------
% Abstract -- JMD requires 150--200 words, Latin characters only (no math),
% structured as background, approach, results, conclusions.
% -----------------------------------------------------------------------------
\begin{abstract}
\todo{Write a 150--200 word abstract covering background, approach, results,
and conclusions. Lorem ipsum follows as a length placeholder.}
\lipsum[1][1-9]
\end{abstract}

\maketitle

% =============================================================================
% Body -- IMRaD structure per JMD submission instructions
% =============================================================================
\section{Introduction}
\label{sec:introduction}
\lipsum[1-3]

\section{Background}
\label{sec:background}
\subsection{Tensegrity Structures}
\lipsum[4]

\subsection{Multifidelity Bayesian Optimization}
\lipsum[5]

\section{Methods}
\label{sec:methods}
\subsection{Design Parameterization and Fabrication}
\lipsum[6]

\subsection{Simulation and Experimental Characterization}
\lipsum[7]

\subsection{Multifidelity Bayesian Optimization Framework}
\lipsum[8]

\section{Results}
\label{sec:results}
\lipsum[9-10]

\section{Discussion}
\label{sec:discussion}
\lipsum[11]

\section{Conclusions}
\label{sec:conclusions}
\lipsum[12]

% =============================================================================
% Back matter
% =============================================================================
\section*{Acknowledgment} % ASME requests this exact spelling, singular.
\todo{Acknowledge funding sources (e.g., BYU MRG, NASA Utah Space Grant
Consortium), facilities, and individual contributors.}

\section*{Funding Data}
\begin{itemize}
  \item \todo{Grant name and number (per ASME Funding Data requirements).}
\end{itemize}

\section*{Conflict of Interest}
The authors declare no conflict of interest.

% -----------------------------------------------------------------------------
% References (ASME numeric style; uses asmejour.bst)
% -----------------------------------------------------------------------------
\bibliographystyle{asmejour}
\bibliography{references}

\end{document}
